\documentclass[aps,pra,reprint,superscriptaddress,nofootinbib,longbibliography]{revtex4-2}

\usepackage{amsmath}
\usepackage{amssymb}
\usepackage{amsthm}
\usepackage{bm}
\usepackage{graphicx}
\usepackage{hyperref}
\usepackage{orcidlink}

\newcommand{\bra}[1]{\langle #1 |}
\newcommand{\ket}[1]{| #1 \rangle}
\newcommand{\braket}[2]{\langle #1 | #2 \rangle}
\newcommand{\Tr}{\mathrm{Tr}}
\newcommand{\eps}{\varepsilon}
\newcommand{\HH}{\mathcal{H}}
\newcommand{\EE}{\mathcal{E}}
\newcommand{\SSys}{\mathcal{S}}
\newcommand{\E}{\mathbb{E}}
\newcommand{\PP}{\mathbb{P}}
\newcommand{\AAlg}{\mathcal{A}}
\newcommand{\MM}{\mathcal{M}}

\begin{document}

\title{Quasi-Orthogonality in High-Dimensional Hilbert Spaces:\\
A Geometric Companion to Environmental Decoherence}

\author{Karl Svozil\,\orcidlink{0000-0001-6554-2802}}
\affiliation{Institute for Theoretical Physics, TU Wien,
Wiedner Hauptstrasse 8-10/136, A-1040 Vienna, Austria}
\email{karl.svozil@tuwien.ac.at}

\date{\today}

\begin{abstract}
The emergence of classicality from unitary quantum mechanics requires the suppression of macroscopic interference and is conventionally modeled dynamically through environment-induced decoherence. We argue that a substantial portion of this story is \emph{kinematic}: the geometry of a high-dimensional Hilbert space alone forces generic vectors to be approximately orthogonal. Combining the exponential packing bound for $\eps$-quasi-orthogonal sets with L\'evy's concentration of measure on the unit sphere, we show that an environment of $N$ particles supports a doubly-exponential supply of $\eps$-quasi-orthogonal pure states, and that two Haar-random vectors have expected modulus-squared overlap $1/d$ with sub-Gaussian tails. We make the kinematic and dynamical inputs explicit: the geometric argument supplies an enormous reservoir of nearly non-interfering branches, but a genuine dynamical assumption---effective pair-typicality on the relevant subspace, generated by post-einselection conditional Hamiltonians and fast-scrambling chaotic dynamics---is still required to guarantee that physical environmental histories actually populate this reservoir. We then connect the framework to the infinite-tensor-product programme of von Neumann, Hepp, Grangier--Van Den Bossche and Svozil, showing that quasi-orthogonality at finite $N$ is the operational shadow of strict sectorization at $N=\infty$ and presages the transition from Type~I to hyperfinite Type~II$_1$ structure. Self-contained proofs of the geometric statements are given in the appendices.
\end{abstract}

\maketitle

\section{Introduction}

The quantum measurement problem asks how the continuous and reversible Schr\"odinger evolution gives rise to the discrete, irreversible classical world we observe~\cite{schlosshauer-2007,schlosshauer-2019}. Operationally, the question reduces to why the off-diagonal interference terms of a system's reduced density matrix are unobservable on macroscopic scales.

The urgency of this question was forcefully articulated by Schr\"odinger in a 1952 colloquium~\cite{schroedinger-interpretation}. He noted that if the wave equation is universally valid, then alternative macroscopic outcomes must all ``really happen simultaneously.'' Schr\"odinger argued that if nature actually took this form, we should find our surroundings rapidly blurring together into a ``quagmire, or sort of a featureless jelly''---a variant of his cat metaphor. For example, a watch left unobserved in a drawer overnight should dissolve into a smeared probability distribution. This ``jellification'' argument poses a sharp, physical question: if unitary quantum mechanics holds at all scales, what fundamental mechanism prevents the simultaneous alternatives of the macroscopic world from blurring into an unobservable, featureless smear?

Bell coined ``FAPP'' (\emph{for all practical purposes}) to express skepticism toward explanations that rest on the practical impossibility of observing macroscopic interference rather than on a fundamental physical mechanism~\cite{bell-a}. Modern decoherence theory~\cite{joos1985emergence,RevModPhys.75.715,joos2003decoherence,schlosshauer-2007} answers this by showing that interaction with a macroscopic environment dynamically drives the off-diagonal terms toward zero in a preferred (pointer) basis selected by the system--environment coupling.

In this Letter we examine a complementary, purely \emph{geometric} contribution to the same suppression. Inspired by recent work in machine learning on the storage of features in high-dimensional embeddings~\cite{elhage2022superposition,bricken2023monosemanticity}, where relaxing exact orthogonality to $\eps$-quasi-orthogonality unlocks an exponential capacity for distinct directions, we transport the same elementary geometry into the Hilbert space of a macroscopic environment. The relevant tools---the exponential packing of quasi-orthogonal unit vectors~\cite{kainen1993quasiorthogonal,tao2012topics,vershynin2018high} and the concentration of measure on high-dimensional spheres~\cite{ledoux2001concentration,milman1986asymptotic}---are classical, but their direct foundational application to the measurement problem is, to our knowledge, underemphasized in the literature.

Our central claim is moderate. Geometry alone guarantees a vast reservoir of nearly non-interfering pure states; this is a kinematic fact about $\HH$ that holds independently of any Hamiltonian. Whether \emph{physical} environmental branches occupy this reservoir is a dynamical question, addressed by the eigenstate thermalization hypothesis (ETH)~\cite{deutsch1991quantum,srednicki1994chaos,rigol2008thermalization,dalessio2016quantum}, by results on scrambling and typicality of unitaries generated by chaotic dynamics~\cite{hayden2007black,brown2018second,maldacena2016bound}, and by the theory of approximate unitary $k$-designs~\cite{brandao2016local,harrow2009random}. We make this division of labor explicit, and we deliberately resist the temptation---present in some informal accounts---to elevate the geometric observation into a stand-alone resolution of the measurement problem. The geometric argument quantitatively defuses Schr\"odinger's ``jellification'' worry; it does not, by itself, dispose of Bell's deeper objection that FAPP suppression is not exact suppression.

A further conceptual benefit, developed in Sec.~V, is that the framework provides an operational, finite-$N$ language for the structural mechanisms long studied in the algebraic and infinite-tensor-product traditions of von Neumann~\cite{vonNeumann1939}, Murray--von Neumann~\cite{Murray1936}, Hepp~\cite{hepp-1972}, Grangier--Van Den Bossche~\cite{grangier2005,van-den-bossche-2023-a}, and Svozil~\cite{svozil-2025-u}. What appears as strict sectorization or a Type-I-to-Type-II$_1$ transition in the $N=\infty$ limit appears in our framework as exponential quasi-orthogonal packing at finite (but very large) $N$.

\section{Quasi-Orthogonality and Exponential Packing}

In a $d$-dimensional Hilbert space $\HH_d$, the maximum number of mutually orthogonal unit vectors is exactly $d$. Two unit vectors $\ket{\psi_i},\ket{\psi_j}$ are called \emph{$\eps$-quasi-orthogonal} if
\begin{equation}
    |\braket{\psi_i}{\psi_j}|^2 \le \eps,
\end{equation}
and a set $\{\ket{\psi_i}\}_{i=1}^{N}$ is $\eps$-quasi-orthogonal if every pair satisfies this bound. The off-diagonal terms in a reduced density matrix arising from such environmental records are weighted by $\braket{E_j}{E_i}$ and are therefore suppressed in modulus by at most $\sqrt{\eps}$.

The cardinality $N_{\eps}(d)$ of the largest such set scales not as $d$ (the orthogonal-set baseline) but exponentially in $d$ once $\eps$ is held fixed and positive:
\begin{equation}
    \exp \bigl(c\,\eps\, d\bigr) \;\le\; N_{\eps}(d) \;\le\; \exp \bigl(C\,\eps\, d\bigr),
    \label{eq:JL}
\end{equation}
for absolute constants $c,C>0$, valid in the regime $2/d\le\eps\le 1/2$ (the trivial orthonormal case $\eps=0$ saturates at $N_0(d)=d$ and is not part of this regime). The lower bound follows from a probabilistic construction (Appendix~\ref{app:packing}); the upper bound is a Kabatiansky--Levenshtein-type estimate~\cite{kabatiansky1978bounds,conway1999sphere,alon2008problems}. The linear-in-$\eps$ scaling in the exponent reflects the squared-overlap convention used in the definition above: an $\eps$-quasi-orthogonal set in this convention is a code of minimum angular distance $\arccos\sqrt{\eps}$, whose spherical-cap volume scales as $(1-\eps)^{\Theta(d)}$. Had we instead defined quasi-orthogonality through the amplitude threshold $|\braket{\psi_i}{\psi_j}|\le\eta$, the corresponding bound would read $\exp(c'\eta^2 d)$, recovering the form familiar from the Johnson--Lindenstrauss flattening lemma~\cite{johnson-1984,alon2016optimal,vershynin2018high}; the two conventions are related by $\eps=\eta^2$. Both are manifestations of concentration of measure in high dimensions.

Equation~\eqref{eq:JL} has a striking implication for many-body physics. For a quantum system of $N$ qubits, $d=2^N$, and the number of available $\eps$-quasi-orthogonal directions scales as
\begin{equation}
    N_{\eps}(2^N) \;\sim\; \exp \bigl(c\,\eps\,2^N\bigr),
\end{equation}
i.e.\ \emph{doubly} exponential in the particle number. For a macroscopic environment with $N \sim 10^{23}$, this number exceeds any conceivable bound on the count of distinct macroscopic histories in the observable universe, by an astronomical margin.

\section{L\'evy's Lemma and Typical Overlaps}

The packing bound \eqref{eq:JL} concerns extremal configurations. The complementary statement, that \emph{generic} (Haar-random) pairs of pure states are already quasi-orthogonal, is a consequence of L\'evy's lemma~\cite{ledoux2001concentration,milman1986asymptotic}.

Identify $\HH_d \cong \mathbb{C}^d$ with $\mathbb{R}^{2d}$ and let $\mathbb{S}^{2d-1}$ denote the unit sphere equipped with the uniform (Haar) measure. For an $L$-Lipschitz function $f:\mathbb{S}^{2d-1}\to\mathbb{R}$, in the Gromov--Milman normalization~\cite[Ch.~1]{ledoux2001concentration},
\begin{equation}
    \PP \Bigl(\bigl|f(\ket{\psi})-\E f\bigr|\ge \delta\Bigr)
    \;\le\; 2\exp \Bigl[-\frac{(2d-1)\,\delta^2}{9\pi^3 L^2}\Bigr].
    \label{eq:Levy}
\end{equation}
(Other normalizations in the literature differ by a constant factor inside the exponent; only the qualitative dimension dependence matters here.) Take $f(\ket{\psi})=|\braket{\phi}{\psi}|^2$ for a fixed reference state $\ket{\phi}$. Two short calculations (Appendix~\ref{app:overlap}) yield
\begin{equation}
    \E_{\ket{\psi}\sim\mu_{\rm Haar}} \bigl[\,|\braket{\phi}{\psi}|^2\,\bigr] \;=\; \frac{1}{d},
    \qquad L \le 2.
    \label{eq:meanLip}
\end{equation}
Substituting into~\eqref{eq:Levy} gives the sub-Gaussian tail bound
\begin{equation}
    \PP \Bigl(|\braket{\phi}{\psi}|^2 \ge \tfrac{1}{d}+\delta\Bigr)
    \;\le\; 2\exp \Bigl[-\frac{(2d-1)\,\delta^2}{36\pi^3}\Bigr].
    \label{eq:overlapTail}
\end{equation}

For this particular observable a sharper and fully explicit statement is available. Standard unitary-invariance arguments show that for fixed $\ket{\phi}$ and Haar-distributed $\ket{\psi}\in\mathbb{S}^{2d-1}$ the random variable $X=|\braket{\phi}{\psi}|^{2}$ has the $\mathrm{Beta}(1,d{-}1)$ distribution, with explicit survival function
\begin{equation}
    \PP\bigl(\,|\braket{\phi}{\psi}|^{2}\ge \eps\,\bigr) \;=\; (1-\eps)^{d-1}
    \;\le\; e^{-(d-1)\eps}.
    \label{eq:overlapExact}
\end{equation}

\begin{figure}[ht]
\centering
\includegraphics[width=\columnwidth]{overlap_survival_aps_monochrome.pdf}
\caption{Exact Haar survival probability for
$X=|\langle \phi|\psi\rangle|^2$ in $\mathbb{C}^d$:
$\mathbb{P}(X\geq\varepsilon)=(1-\varepsilon)^{d-1}$.
Curves are shown for $d=4,8,16,32$.
The rapid collapse of the right tail illustrates the finite-size onset of quasi-orthogonality.}
\label{fig:overlap-survival}
\end{figure}

To make the overlap statistics concrete, consider the smallest nontrivial many-body environment, $N=3$ qubits, so that $d=2^N=8$. Fix $\ket{\phi}=\ket{000}$ and draw $\ket{\psi}$ Haar-randomly from $\mathbb{S}^{15}$. Then
\[
X:=|\braket{\phi}{\psi}|^2
\]
has the exact Beta$(1,7)$ density
\[
p_8(x)=7(1-x)^6,\qquad 0\le x\le 1,
\]
with mean $\E[X]=1/8$ and median
\[
x_{1/2}=1-2^{-1/7}\approx 9.4\times10^{-2}.
\]
Even at this modest dimension, large overlaps are already suppressed:
\begin{align}
&\PP(X\ge 1/4)=0.75^7\approx 1.34\times10^{-1},      \nonumber
\\
&\PP(X\ge 1/2)=2^{-7}\approx 7.8\times10^{-3}.     \nonumber
\end{align}
For $N=4$ qubits ($d=16$), the corresponding probabilities become
\begin{align}
&\PP(X\ge 1/4)=0.75^{15}\approx 1.33\times10^{-2},  \nonumber
\\
&\PP(X\ge 1/2)=2^{-15}\approx 3.1\times10^{-5}.   \nonumber
\end{align}
Fig.~\ref{fig:overlap-survival} displays the corresponding survival functions for several values of $d$ and makes visually clear how rapidly the right tail of the overlap distribution is pushed toward zero as $d$ increases.

The exponent in~\eqref{eq:overlapExact} is linear in $\eps$, whereas~\eqref{eq:overlapTail} (a Lipschitz consequence of L\'evy's lemma) yields only a quadratic-in-$\delta$ Gaussian envelope. For the present application either bound suffices; it is~\eqref{eq:overlapExact}, however, that is responsible for the linear-in-$\eps$ scaling of the packing exponent in~\eqref{eq:JL}, and we use it for the proof of that bound in Appendix~\ref{app:packing}. For $d=2^{N}$ with $N$ macroscopic, the right-hand side of~\eqref{eq:overlapExact} is doubly exponentially small in $N$ for any fixed $\eps>0$, decaying as $\exp[-(2^N-1)\eps]$: \emph{any two Haar-random pure states in a macroscopic Hilbert space are $\eps$-quasi-orthogonal with overwhelming probability}.

\section{Application to the Measurement Process}

Consider a system $\SSys$ coupled to an apparatus--environment $\EE$ initialised in $\ket{E_0}$:
\begin{equation}
    \ket{\Psi(0)} = \Bigl(\sum_i c_i \ket{s_i}\Bigr)\otimes\ket{E_0}.
\end{equation}
A measurement-like unitary correlates pointer outcomes with environmental records:
\begin{equation}
    \ket{\Psi(t)} = \sum_i c_i \ket{s_i}\otimes\ket{E_i(t)}.
\end{equation}
The reduced density matrix on $\SSys$ is
\begin{equation}
    \rho_{\SSys}(t) = \sum_{i,j} c_i c_j^{*}\,\braket{E_j(t)}{E_i(t)}\,\ket{s_i} \bra{s_j}.
\end{equation}
Off-diagonal coherences are weighted by the environmental overlaps $\braket{E_j(t)}{E_i(t)}$.

\subsection*{The kinematic claim.}
Throughout this section we adopt natural units $k_B=\hbar=1$, so that entropies are dimensionless and inverse temperatures have units of time. \emph{If} the branch states $\{\ket{E_i(t)}\}_{i\ne j}$ are well-modeled as a typical pair on the relevant high-dimensional subspace of $\HH_{\EE}$ (in a sense made precise below), then
\begin{equation}
    \E\bigl[\,|\braket{E_j(t)}{E_i(t)}|^2\,\bigr] \;=\; \frac{1}{d_{\rm eff}},
\end{equation}
with $d_{\rm eff}$ the dimension of the subspace explored by the dynamics. For a generic chaotic many-body environment the accessible subspace is the microcanonical shell, so
\begin{equation}
    d_{\rm eff} \;\sim\; e^{S},
    \qquad S \sim N,
    \label{eq:dEntropy}
\end{equation}
with $S$ the (dimensionless) microcanonical entropy of the relevant environmental sector and $N$ the number of constituent degrees of freedom; in conventional units the right-hand sides of~\eqref{eq:dEntropy} read $e^{S/k_B}$ and $Nk_B$, respectively (see also Appendix~\ref{app:deff}). Equation~\eqref{eq:dEntropy} ties the geometric parameter $d_{\rm eff}$ directly to standard thermodynamics: the mean-square overlap is exponential in the entropy of the environment, so the typical overlap amplitude is of order $e^{-S/2}$. This is vastly stronger than needed to render coherences operationally invisible on any timescale shorter than the Poincar\'e recurrence time, which for a generic chaotic many-body system is astronomically large---heuristic arguments suggest a scaling at least as $\exp(d_{\rm eff})\sim\exp(e^{S})$, but the precise value is model-dependent and irrelevant to our argument~\cite{peres1982recurrence,wallace2015recurrence}.

\subsection*{The dynamical assumption.}
The conditional ``if'' above is doing real work and must not be elided. The branch states $\ket{E_i(t)}$ are not random vectors plucked from $\HH_{\EE}$; they are deterministic images of $\ket{E_0}$ under specific unitaries $U_i(t)$ generated by the system--environment Hamiltonian.

A natural mechanism, however, makes the typicality model defensible. Once the einselection mechanism of decoherence theory~\cite{zurek1981pointer,RevModPhys.75.715} has singled out a pointer basis $\{\ket{s_i}\}$, the system--environment coupling reduces, in this basis, to the schematic form
\begin{equation}
    H_{\rm int} \;=\; \sum_i \ket{s_i} \bra{s_i} \otimes H_{\EE,i},
    \label{eq:condH}
\end{equation}
i.e.\ the environment evolves under a \emph{different conditional Hamiltonian} $H_{\EE,i}$ for each pointer index. This post-einselection form is not generic at the level of Hamiltonians on $\HH_\SSys\otimes\HH_\EE$---a generic interaction reads $\sum_a A_a\otimes B_a$---but is the form to which generic couplings reduce once the pointer basis has been selected. The geometric argument below is silent on this selection and presupposes it (cf.\ the preferred-basis caveat at the end of this section). When the conditional generators are non-commuting and chaotic---as they generically are for macroscopic environments with local interactions---the unitaries $U_i(t)=e^{-iH_{\EE,i}t}$ rapidly delocalise the common initial state $\ket{E_0}$ across the microcanonical shell. A standard benchmark for this delocalisation is the fast-scrambling time~\cite{hayden2007black,maldacena2016bound,brown2018second},
\begin{equation}
    t_{*} \;\gtrsim\; \frac{\beta}{2\pi}\,\log N,
    \label{eq:scramble}
\end{equation}
where $\beta=1/T$ in the natural units adopted above and $N$ is the number of local degrees of freedom. The qualifier ``fast'' refers precisely to this \emph{logarithmic} dependence on system size: it is to be contrasted with the much slower power-law scrambling characteristic of diffusive or integrable dynamics, in which $t_*$ grows polynomially in $N$. The bound is saturated, in particular, by black-hole-like and matrix-model dynamics~\cite{hayden2007black,maldacena2016bound,brown2018second}.

For $t\gtrsim t_*$, each branch state $\ket{E_i(t)}$ is effectively delocalised across the accessible microcanonical shell, and one expects the \emph{pair} $\{\ket{E_i(t)},\ket{E_j(t)}\}$ with $i\ne j$ to behave, at the level of low-degree polynomial observables such as the squared overlap, like a typical pair of points on that shell. The relevant technical input is the theory of approximate unitary $k$-designs: the chaotic conditional dynamics generated by~\eqref{eq:condH} can be argued to form an approximate $2$-design on the accessible shell on a time of order~\eqref{eq:scramble}~\cite{brandao2016local,harrow2009random,hayden2007black}, which is exactly the property required for the second moment of the overlap (and hence its mean-square) to match the Haar value $1/d_{\rm eff}$. We emphasize, however, that this is a \emph{modeling} assumption, not a theorem: the two branch states are deterministic images of a common initial state under different unitaries and are not literally independent Haar-random vectors. The defensible statement is the weaker one that the pair-distribution induced by the ensemble of conditional dynamics, varied over the natural chaotic class to which $H_{\EE,i}$ belongs, approaches the Haar-pair distribution at the level of polynomial observables, with corrections subleading in $1/d_{\rm eff}$ and controlled by the $2$-design error. Under this (standard but nontrivial) assumption, the geometric estimates of Sec.~III apply.

The supporting dynamical results---ETH~\cite{deutsch1991quantum,srednicki1994chaos,rigol2008thermalization,dalessio2016quantum}, Page's calculation~\cite{page1993average}, Reimann's typicality theorems~\cite{reimann2008foundation,linden2009quantum}, and approximate-design results~\cite{brandao2016local,harrow2009random}---all converge on the same picture. For integrable, many-body-localised, or weakly-coupled systems these assumptions can fail, and the geometric reservoir, while available, is not occupied by physical histories. The geometric statement is therefore a kinematic \emph{license}, conditional on a dynamical input that is generically but not universally satisfied.

\subsection*{What this does not solve.}
Two further caveats are essential and are sometimes blurred in informal presentations.

\paragraph{The preferred-basis problem.} Geometric quasi-orthogonality is basis-independent: if $\{\ket{E_i}\}$ are mutually quasi-orthogonal, so is any unitary rotation of them. Decoherence theory's einselection mechanism~\cite{zurek1981pointer,RevModPhys.75.715} singles out a \emph{specific} basis (typically quasi-local in position) determined by the form of $H_{\rm int}$ in Eq.~\eqref{eq:condH}. The geometric framework is silent on this selection; it tells us the reservoir is large but not which directions in it are physically realized.

\paragraph{The reality of unselected branches.} As Bell stressed, suppression of off-diagonal terms to any nonzero level---no matter how small---does not by itself convert an improper mixture into a proper one~\cite{bell-a,wallace2012emergent}. To deny ontological status to non-actualized branches one needs an interpretive commitment (collapse, hidden variables, Everettian branching), not a sharper bound on overlaps. Our geometry sharpens the FAPP suppression dramatically, but FAPP, however small, remains FAPP.

\subsection*{Defusing Schr\"odinger's ``jellification'' worry.}
\label{sec:jelly}
Schr\"odinger's fear of a ``featureless jelly''~\cite{schroedinger-interpretation} corresponds mathematically to macroscopic off-diagonal interference terms remaining unsuppressed. The geometric framework provides a quantitative rebuttal, conditional on the dynamical assumptions above. Once the conditional, chaotic dynamics generated by~\eqref{eq:condH} drive the environmental branches into typical regions of the accessible microcanonical shell on the timescale~\eqref{eq:scramble}, the overlap between branches is concentrated near $1/d_{\rm eff}\sim e^{-S}$, with typical amplitude of order $e^{-S/2}$. The cross-terms required to create Schr\"odinger's quagmire are thus mathematically suppressed to a level FAPP-indistinguishable from zero. Moreover, the doubly-exponential packing bound $N_{\eps}(2^N) \sim \exp (c\,\eps\,2^N)$ demonstrates that the macroscopic Hilbert space has vastly more than enough ``room'' to host simultaneous alternatives without them kinematically interfering. The alternatives do not blur into a single smeared distribution; rather, because they are driven into approximately orthogonal quasi-sectors, the macroscopic contours within each sector remain sharply defined---to the precision controlled by $e^{-S/2}$. The geometric argument thus quantifies the FAPP cushion that prevents jellification, without claiming to resolve the deeper interpretive questions that remain.

\section{Connections to the Literature}

\subsection*{A. Canonical typicality.}
The framework here is a near-relative of canonical typicality~\cite{popescu2006entanglement,goldstein2006canonical,reimann2008foundation,linden2009quantum,gemmer2009quantum}: L\'evy's lemma applied to the partial trace shows that almost every pure state of a large composite system has a reduced state close to the canonical thermal state on a small subsystem. Our application is dual, asking instead about \emph{pairwise} overlaps between distinct macroscopic records.

\subsection*{B. Everettian branching.}
The doubly-exponential packing capacity makes precise an intuition often invoked in defenses of the Many-Worlds Interpretation~\cite{wallace2012emergent,saunders2010many}: the Hilbert space ``has room'' for an unfathomable number of branches without their kinematically interfering. We emphasize, however, that kinematic room is necessary but not sufficient for the Everettian picture, which additionally requires a dynamical branching structure and a probabilistic interpretation~\cite{wallace2012emergent}.

\subsection*{C. Quantum Darwinism.}
Zurek's quantum Darwinism~\cite{zurek2009quantum,zwolak2017redundancy} requires the environment to redundantly encode a single classical record across many fragments. The geometric capacity demonstrated here underwrites the very possibility of such redundant encoding: the environment is large enough to host many nearly independent copies, by the same packing bound~\eqref{eq:JL}.

\subsection*{D. Random-matrix and chaotic-dynamics approaches.}
Recent work has connected typicality to dynamical chaos through random-matrix models of quantum thermalization~\cite{dalessio2016quantum,cotler2019chaos}, and through bounds on scrambling rates~\cite{maldacena2016bound}. Our geometric perspective is consistent with these and clarifies why such approaches succeed: once dynamics drives a state into a typical region of a high-dimensional subspace, the kinematic conclusions of Sec.~III take over.

\subsection*{E. Sectorization, factorization, and finite-$N$ shadows.}
\label{sec:sectors}
A complementary---and historically deeper---route to FAPP irreversibility proceeds not through finite-dimensional concentration of measure, but through the structural pathology of \emph{infinite} tensor products. Following von Neumann's analysis of ``incomplete infinite direct products''~\cite{vonNeumann1939} and its development in algebraic quantum field theory~\cite{haag-1996,doplicher-1969,doplicher-1971} and in the statistical-mechanical theory of measurement~\cite{hepp-1972,bub-1988}, the infinite tensor product space $\bigotimes_{n=1}^{\infty}\HH_n$ decomposes into mutually orthogonal \emph{sectors}: equivalence classes of sequences $\{\ket{\psi_n}\},\{\ket{\phi_n}\}$ for which
\begin{equation}
    \sum_{n=1}^{\infty}\bigl(1-|\braket{\psi_n}{\phi_n}|\bigr) \;<\; \infty.
    \label{eq:sectorCondition}
\end{equation}
Equation~\eqref{eq:sectorCondition} is the phase-insensitive form appropriate to macroscopic/operational sectors; von Neumann's original 1939 condition for strict vector equivalence retains phase information and reads $\sum_n|1-\braket{\psi_n}{\phi_n}|<\infty$~\cite{vonNeumann1939}, the two coinciding once one quotients by global phases of the individual factors---which is the appropriate identification for sectors interpreted as classical pointer outcomes. Vectors in different sectors are \emph{strictly} orthogonal, and no convergent infinite tensor product of unitaries can interconvert them. Sectorization has been advocated as a structural mechanism for the emergence of classical pointer states by Hepp~\cite{hepp-1972} and, more recently, by Grangier and Van Den Bossche~\cite{grangier2005,van-den-bossche-2023-a}; an extended argument tying nested Wigner's-friend chains to the same limit appears in Svozil~\cite{svozil-2025-u}.

The geometric perspective developed here can be read as the \emph{finite-$N$ shadow} of this picture. Three correspondences make this precise.

\paragraph{(i) Quasi-sectors at finite $N$.}
At finite $N$, the strict equivalence relation defining sectors is replaced by an approximate one at threshold $\eps$: a \emph{quasi-sector} is a maximal $\eps$-quasi-orthogonal family in $\HH_{2^N}$. The packing bound~\eqref{eq:JL} yields $N_{\eps}(2^N)\sim e^{c\eps\,2^N}$ such quasi-sectors. This finite-$N$ notion is the shadow of the infinite-tensor-product criterion: for product vectors $\ket{\Psi}=\bigotimes_n\ket{\psi_n}$ and $\ket{\Phi}=\bigotimes_n\ket{\phi_n}$, the condition~\eqref{eq:sectorCondition} means that the infinite product of moduli $\prod_n|\braket{\psi_n}{\phi_n}|$ converges to a nonzero limit and the two vectors belong to the same sector, whereas divergence of the sum forces the product overlap to vanish and places the vectors in different, mutually orthogonal sectors. In this sense the finite-$N$ packing bound is an operational version of sectorization: small per-component mismatches can accumulate into a macroscopic suppression of the total overlap. Equivalently, the central inequality used by Svozil~\cite{svozil-2025-u} and by Van Den Bossche--Grangier~\cite{van-den-bossche-2023-a}---namely, $\prod_{n=1}^N |\braket{\psi_n}{\phi_n}| \le \exp\bigl[-\sum_{n=1}^N (1-|\braket{\psi_n}{\phi_n}|)\bigr]$, so that any persistent positive lower bound on per-site mismatch drives the total overlap to zero exponentially in $N$---is, at finite $N$, the multiplicative form of the quasi-orthogonality bound and is controlled quantitatively by Eqs.~\eqref{eq:overlapTail} and~\eqref{eq:overlapExact}.

\paragraph{(ii) From Type~I to hyperfinite Type~II$_1$.}
The local algebra at finite $N$ is the Type~I$_{2^N}$ factor $\AAlg_N=\MM_{2^N}(\mathbb{C})$, whose minimal projections are rank one and on which the normalized trace takes the discrete set of values $\{k\,2^{-N}: 0\le k\le 2^N\}$. The inductive limit of the algebras $\AAlg_N$ along the inclusions $a\mapsto a\otimes\mathbf{1}$ is the uniformly hyperfinite (UHF) $C^*$-algebra of type $2^\infty$; in the Gel'fand--Naimark--Segal (GNS) representation of the unique tracial state compatible with these inclusions, its weak operator closure is the hyperfinite Type~II$_1$ factor $\mathcal{R}$~\cite{Murray1936,Jones-vNA}, which is \emph{diffuse}: it has no minimal projections, and the trace takes a continuum of values densely in $[0,1]$. The doubly-exponential proliferation of $\eps$-quasi-orthogonal projections derived in Sec.~II is the kinematic precursor of this diffuseness in two senses, which can be made quantitative.

First, the discrete set of trace values $\{k\,2^{-N}\}$ becomes increasingly fine as $N$ grows, and indeed dense in $[0,1]$ in the limit; this is the trivial statement underlying the well-known fact that the limit factor admits a continuous trace. Second, and more substantively, the number of mutually $\eps$-quasi-orthogonal rank-one projections at level $N$---each of trace $2^{-N}$---is $N_\eps(2^N)\sim e^{c\eps\cdot 2^N}$, so the multiplicity of operationally indistinguishable projections at each trace value grows doubly-exponentially. Under the equivalence furnished by quasi-orthogonality at threshold $\eps$, no single rank-one projection is operationally distinguishable from a family of $\sim e^{c\eps\cdot 2^N}$ near-equivalents, and the cardinality of equivalence classes of rank-one projections compatible with the trace-value $2^{-N}$ diverges in the limit. Diffuseness of $\mathcal{R}$ does not appear suddenly at $N=\infty$; it is approached gradually, with the geometric capacity providing a sharp finite-$N$ measure of proximity to the continuum limit.

For the corresponding Type~III structures encountered in algebraic quantum field theory and in recent gravitational settings~\cite{penington2023algebras}, an analogous shadow exists but is more subtle: Type~III factors lack any normalized trace, so the appropriate finite-$N$ proxy is not a projection count but the entanglement-spectrum statistics, which our framework also constrains via Eq.~\eqref{eq:dEntropy}.

\paragraph{(iii) Two routes to FAPP, one mathematical fact.}
Svozil~\cite{svozil-2025-u} and Grangier--Van Den Bossche~\cite{grangier2005,van-den-bossche-2023-a} obtain irreversibility by passing to the strict $N=\infty$ limit and observing that the resulting transitions cease to be unitary in the sense of the original finite-dimensional theory. Our route obtains operational irreversibility at finite (but enormous) $N$, where unitarity remains mathematically intact but the mean-square size of off-diagonal coherences is suppressed by $\sim e^{-S}$, with the typical amplitude suppressed by $\sim e^{-S/2}$, where $S\sim N$ in natural units (and $S\sim Nk_B$ in conventional ones). The two routes are duals of one another: the infinite-tensor-product argument trades finite control for a clean limit; the geometric argument trades a clean limit for finite control. Both rest on the same mathematical phenomenon---the orthogonality (or near-orthogonality) of generic vectors in product spaces of high cardinality---and both remain vulnerable to the same Bell-style objection that ``FAPP is not exact.'' Whether one prefers the kinematic shadow or the limit it shadows is, in the end, partly a question of one's willingness to accept transfinite mathematical objects as physically meaningful.

\subsection*{F. Machine-learning embeddings.}
The motivating analogy to high-dimensional vector embeddings in language models and superposition codes~\cite{elhage2022superposition,bricken2023monosemanticity} is more than rhetorical: the same packing bound that enables a network with $d$ hidden dimensions to store $\sim e^{c\eps\, d}$ approximately-disentangled features (in the squared-correlation convention adopted here, equivalent to $e^{c'\eta^2 d}$ in the more common amplitude-correlation convention with $\eta=\sqrt{\eps}$) enables a Hilbert space of dimension $d$ to host that many approximately-classical records. Both phenomena are facets of high-dimensional geometry and concentration of measure.

\section{Discussion}

The geometric perspective supplies a quantitative kinematic floor on environmental orthogonality but does not, by itself, vindicate Bell's worry about FAPP. What it does accomplish is to disentangle two distinct contributions to decoherence that are often bundled together. The first is dynamical: the system--environment coupling \eqref{eq:condH} generates entanglement, selects a quasi-local pointer basis, and drives initial overlaps to small values on the fast-scrambling timescale~\eqref{eq:scramble}. The second is geometric: once states reach a typical region of a high-dimensional space, their mean-square overlaps are concentrated near $1/d_{\rm eff}\sim e^{-S}$, so the typical overlap amplitude is of order $e^{-S/2}$, regardless of how the states were generated. The two contributions are logically independent, and our analysis is a clean separation of variables.

A natural follow-up is to ask whether the geometric reservoir can be put under direct experimental pressure. Mesoscopic interference experiments---in superconducting qubits, trapped ions, or matter-wave interferometers~\cite{arndt2014testing}---probe regimes in which $d_{\rm eff}$ is moderate and the typicality assumption is partially violated; deviations from the predicted overlap distribution, in particular slower-than-$e^{-S}$ decay of coherences, would directly constrain how rapidly chaotic dynamics drives states into typicality, and would diagnose obstructions such as approximate conservation laws, weak ergodicity breaking, or many-body localization. Such measurements---probing the ``approach to typicality'' rather than its endpoint---are within reach of current platforms with $N\sim 10$--$100$ controllable degrees of freedom, where $d_{\rm eff}$ is large enough for the geometric statements to bite but small enough for genuine deviations from Haar statistics to be experimentally resolvable. We leave a detailed proposal to future work.

\section{Conclusion}

We have argued that the suppression of macroscopic interference has a clean kinematic component---the exponential packing and concentration of measure of high-dimensional Hilbert spaces---which, under standard chaotic-dynamics assumptions, quantitatively defuses Schr\"odinger's ``jellification'' worry~\cite{schroedinger-interpretation} and complements the dynamical mechanism of environmental decoherence. In a $d$-dimensional space, an exponential-in-$d$ supply of $\eps$-quasi-orthogonal pure states is available, and Haar-random pairs satisfy $\E|\braket{\phi}{\psi}|^2 = 1/d$ with the explicit Beta tail $\PP(|\braket{\phi}{\psi}|^2\ge\eps)=(1-\eps)^{d-1}$ and the corresponding sub-Gaussian L\'evy bound.

For macroscopic $N$, with $d_{\rm eff}\sim e^S$ and $S$ extensive in $N$ (i.e.\ $S\sim N$ in natural units, or $S\sim Nk_B$ in conventional ones), this makes the mean-square inter-branch overlap exponentially small in $N$. Whenever the conditional dynamics of Eq.~\eqref{eq:condH} drive the environment into a typical region of its microcanonical shell on the fast-scrambling timescale~\eqref{eq:scramble}, the number of available quasi-orthogonal directions is doubly exponential in $N$. This provides the mathematical room for simultaneous branches to co-exist without blurring into a macroscopic quagmire, offering a finite-$N$, operational shadow of the strict sectorization seen in the infinite-tensor-product limit, and a finite-$N$ precursor of the diffuse trace structure of the hyperfinite Type~II$_1$ factor.

While the geometry does not single out a preferred basis or conceptually convert improper mixtures into proper ones---and to that extent does not resolve the interpretive core of the measurement problem---it makes precise, and surprisingly extreme, the size of the FAPP cushion on which the classical world rests. Conditional on the standard chaotic-typicality assumptions, it explains why macroscopic wavefunctions do not manifest as a featureless jelly.

\appendix

\section{Mean and Lipschitz Bound for Overlap}\label{app:overlap}

Throughout this appendix, $\|\cdot\|$ denotes the Euclidean (chordal) norm $\|\ket{\psi}\|=\sqrt{\braket{\psi}{\psi}}$ on $\HH_d\cong\mathbb{C}^d\cong\mathbb{R}^{2d}$, restricted to the unit sphere $\mathbb{S}^{2d-1}$. This is the norm appearing in L\'evy's lemma~\eqref{eq:Levy}, and is to be distinguished from the trace, Hilbert--Schmidt, and operator norms used elsewhere in the paper.

\begin{proof}[Mean]
Let $\ket{\phi}\in\HH_d$ be fixed and let $\ket{\psi}$ be Haar-distributed on the unit sphere. By unitary invariance, $\E|\braket{\phi}{\psi}|^2$ depends only on the dimension. Choose any orthonormal basis $\{\ket{e_k}\}_{k=1}^d$. Then
\[
1=\E\sum_{k=1}^d |\braket{e_k}{\psi}|^2 = d\cdot \E|\braket{e_1}{\psi}|^2,
\]
so $\E|\braket{\phi}{\psi}|^2 = 1/d$.
\end{proof}

\begin{proof}[Lipschitz constant]
Let $f(\ket{\psi})=|\braket{\phi}{\psi}|^2$ on the unit sphere with the chordal Euclidean metric $\|\,\cdot\,\|$ defined above. For unit vectors $\ket{\psi_1},\ket{\psi_2}$,
\begin{align*}
|f(\psi_1)-f(\psi_2)| &= \bigl|\,|\braket{\phi}{\psi_1}|^2 - |\braket{\phi}{\psi_2}|^2 \,\bigr| \\
&\le \bigl(|\braket{\phi}{\psi_1}|+|\braket{\phi}{\psi_2}|\bigr)\bigl|\braket{\phi}{\psi_1}-\braket{\phi}{\psi_2}\bigr|\\
&\le 2\,|\braket{\phi}{\psi_1-\psi_2}| \le 2\,\|\psi_1-\psi_2\|,
\end{align*}
where we used $|\braket{\phi}{\psi_k}|\le 1$ and Cauchy--Schwarz. Hence $L\le 2$.
\end{proof}

\section{Lower Bound on Quasi-Orthogonal Packing}\label{app:packing}

\begin{proof}[Probabilistic construction]
Sample $N$ unit vectors $\ket{\psi_1},\dots,\ket{\psi_N}\in\HH_d$ independently from the Haar measure on $\mathbb{S}^{2d-1}$. For any pair $i\ne j$, by the exact Beta tail~\eqref{eq:overlapExact},
\[
\PP\bigl(|\braket{\psi_i}{\psi_j}|^2 \ge \eps\bigr) \;=\; (1-\eps)^{d-1} \;\le\; e^{-(d-1)\eps}.
\]
By a union bound over the $\binom{N}{2}<N^2/2$ pairs,
\[
\PP\bigl(\exists\, i\ne j:\ |\braket{\psi_i}{\psi_j}|^2 \ge \eps\bigr) \;\le\; \tfrac{1}{2}N^2\,e^{-(d-1)\eps}.
\]
Choosing $N=\bigl\lfloor e^{(d-1)\eps/2-1/2}\bigr\rfloor$ makes this bound less than $1$, so an $\eps$-quasi-orthogonal set of size $N$ exists. Thus $N_{\eps}(d)\ge \exp(c\,\eps\, d)$ with $c=1/2$ for $d$ large.
\end{proof}

The matching upper bound $N_{\eps}(d)\le\exp(C\eps\, d)$ follows from sphere-packing arguments of Kabatiansky--Levenshtein type~\cite{kabatiansky1978bounds,conway1999sphere}: an $\eps$-quasi-orthogonal set is in particular a code of minimum angular distance $\theta=\arccos\sqrt{\eps}$, whose cardinality is bounded above by the inverse of the normalized volume of a spherical cap of half-angle $\theta/2$. The constants $c,C$ are not identical, but both lie in $(0,\infty)$ uniformly in $d$ for $\eps$ in any compact subinterval of $(0,1)$.

We remark that the corresponding bound for the \emph{amplitude} threshold $|\braket{\psi_i}{\psi_j}|\le\eta$ is $\exp(c'\eta^2 d)$ with the well-known Johnson--Lindenstrauss exponent~\cite{johnson-1984,vershynin2018high}; the two are related by $\eps=\eta^2$, which translates the exponent $\eta^2 d$ into $\eps d$ as above.

\section{Effective Dimension and the Microcanonical Entropy}\label{app:deff}

In~\eqref{eq:overlapTail} the relevant dimension is the dimension of the subspace $\HH_{\rm acc}\subset\HH_{\EE}$ \emph{accessible} to the dynamics from the initial state on the timescale considered, not the full Hilbert space dimension. Energy conservation, conserved charges, and finite-time evolution all restrict the accessible subspace. Concretely, for a many-body system with Hamiltonian $H_{\EE}$ and an initial state localized within an energy window $[E,E+\Delta E]$, the accessible subspace is well approximated by the microcanonical shell, and its dimension is
\begin{equation}
    d_{\rm eff} \;=\; \dim \HH_{\rm acc} \;\sim\; e^{S(E)/k_B}
    \label{eq:appS}
\end{equation}
with $S(E)$ the microcanonical entropy. In the natural units adopted in Sec.~IV ($k_B=\hbar=1$), the right-hand side reads $e^{S(E)}$, consistent with~\eqref{eq:dEntropy}. For chaotic many-body dynamics with bounded local interactions, $S(E)$ is extensive in $N$ (Boltzmann scaling), so $d_{\rm eff}$ remains exponential in $N$ and the qualitative conclusions of Sec.~IV survive.

For integrable or many-body-localized systems, conserved local quantities restrict the dynamics to a polynomial-in-$N$ submanifold of the microcanonical shell~\cite{rigol2008thermalization,dalessio2016quantum}, and $d_{\rm eff}$ can be polynomial or even constant in $N$. The geometric reservoir is correspondingly smaller, and the conditional dynamical assumption of Sec.~IV is correspondingly violated. This is a further reason to view the geometric statement as conditional on a dynamical input, and provides a quantitative diagnostic: any experimental observation of slower-than-$e^{-S}$ decoherence is direct evidence of obstructions to effective typicality on the accessible subspace.

\begin{acknowledgments}
This text was partially created and revised with assistance from large language models.
This research was funded in whole or in part by the Austrian Science Fund (FWF) Grant DOI: 10.55776/PIN5424624.
The author acknowledges TU Wien Bibliothek for financial support through its Open Access Funding Programme.
\end{acknowledgments}

\bibliography{svozil}

\end{document}


import numpy as np
import matplotlib.pyplot as plt

# APS-like monochrome style
plt.rcParams.update({
    "font.family": "serif",
    "font.serif": ["Times New Roman", "Times", "DejaVu Serif"],
    "mathtext.fontset": "cm",
    "font.size": 8.5,
    "axes.labelsize": 9,
    "legend.fontsize": 8,
    "xtick.labelsize": 8,
    "ytick.labelsize": 8,
    "axes.linewidth": 0.8,
    "lines.linewidth": 1.7,
    "xtick.direction": "in",
    "ytick.direction": "in",
    "xtick.top": True,
    "ytick.right": True,
    "xtick.major.size": 3.5,
    "ytick.major.size": 3.5,
    "xtick.minor.size": 2.0,
    "ytick.minor.size": 2.0,
    "xtick.major.width": 0.8,
    "ytick.major.width": 0.8,
    "xtick.minor.width": 0.6,
    "ytick.minor.width": 0.6,
    "savefig.bbox": "tight",
    "savefig.pad_inches": 0.02,
    "figure.dpi": 150,
    "savefig.dpi": 300,
})

# Exact Haar survival function:
# P(X >= eps) = (1 - eps)^(d - 1),  X = |<phi|psi>|^2
eps = np.linspace(0.0, 0.85, 600)
dims = [4, 8, 16, 32]
linestyles = ["-", "--", "-.", ":"]
labels = [r"$d=4$", r"$d=8$", r"$d=16$", r"$d=32$"]

fig, ax = plt.subplots(figsize=(3.35, 2.05), constrained_layout=True)

for d, ls, lab in zip(dims, linestyles, labels):
    tail = (1.0 - eps)**(d - 1)
    ax.semilogy(
        eps, tail,
        color="black",
        linestyle=ls,
        label=lab
    )

ax.set_xlabel(r"$\varepsilon$")
ax.set_ylabel(r"$\mathbb{P}(X \geq \varepsilon)$")
ax.set_xlim(0.0, 0.85)
ax.set_ylim(1e-8, 1.2)

ax.minorticks_on()
ax.grid(False)

ax.legend(
    loc="lower left",
    frameon=False,
    handlelength=2.5,
    borderaxespad=0.4
)

for spine in ax.spines.values():
    spine.set_linewidth(0.8)

plt.savefig("overlap_survival_aps_monochrome.pdf")
plt.show()
